\documentclass[11pt,a4paper]{article}

\usepackage[margin=2cm]{geometry}
\usepackage{fontspec}
\usepackage{polyglossia}
\usepackage{amsmath,amssymb,amsthm,mathtools}
\usepackage{booktabs}
\usepackage{array}
\usepackage{enumitem}
\usepackage{xcolor}
\usepackage{hyperref}
\usepackage{titlesec}
\usepackage{fancyhdr}
\usepackage{microtype}
\usepackage{tabularx}

\setmainlanguage{russian}
\setotherlanguage{english}
\setmainfont{Liberation Serif}
\setsansfont{Liberation Sans}
\setmonofont{Liberation Mono}

\definecolor{ruleblue}{HTML}{1F4E79}
\definecolor{rulegray}{HTML}{4A4A4A}
\definecolor{warnred}{HTML}{8B1E1E}
\definecolor{okgreen}{HTML}{1B5E20}
\definecolor{boxbg}{HTML}{F4F7FB}
\definecolor{warnbg}{HTML}{FFF4F0}
\definecolor{okbg}{HTML}{F1F8F2}

\hypersetup{
  colorlinks=true,
  linkcolor=ruleblue,
  urlcolor=ruleblue,
  pdftitle={Робастная устойчивость и критерий Михайлова},
  pdfauthor={MIPT notes}
}

\titleformat{\section}{\large\bfseries\color{ruleblue}}{\thesection.}{0.6em}{}
\titleformat{\subsection}{\normalsize\bfseries\color{ruleblue}}{\thesubsection.}{0.5em}{}
\titlespacing*{\section}{0pt}{1.3em}{0.5em}
\titlespacing*{\subsection}{0pt}{0.9em}{0.35em}

\pagestyle{fancy}
\fancyhf{}
\fancyhead[L]{\small\color{rulegray} Робастная устойчивость}
\fancyhead[R]{\small\color{rulegray} Критерий Михайлова}
\fancyfoot[C]{\thepage}
\renewcommand{\headrulewidth}{0.3pt}

\setlist[enumerate]{leftmargin=1.6em,itemsep=0.25em,topsep=0.35em}
\setlist[itemize]{leftmargin=1.4em,itemsep=0.2em,topsep=0.3em}

\newcommand{\R}{\mathbb{R}}
\newcommand{\jomega}{j\omega}

\newtheorem*{rulebox}{Правило}

\begin{document}

\begin{center}
  {\LARGE\bfseries\color{ruleblue} Робастная устойчивость}\\[0.35em]
  {\Large и критерий Михайлова}\\[0.6em]
  {\large Основные правила решения задач}\\[0.4em]
  {\normalsize\color{rulegray} МФТИ. Интервальные полиномы, годограф Михайлова, теорема Харитонова}
\end{center}

\vspace{0.4em}
\noindent\textcolor{ruleblue}{\rule{\textwidth}{0.8pt}}

\tableofcontents
\vspace{0.3em}
\noindent\textcolor{ruleblue}{\rule{\textwidth}{0.4pt}}

%---------------------------------------------------------------------
\section{Сначала получи характеристический полином}

Для системы с отрицательной обратной связью
\[
W_{\mathrm{замк}}(s)=\frac{W_1(s)}{1+W_1(s)W_2(s)}.
\]
Устойчивость определяется нулями знаменателя:
\[
P(s)=1+W_1(s)W_2(s)=0.
\]

\paragraph{Практический приём.}
Приведи \(W_1\) и \(W_2\) к дробям \(N/D\) и сложи над общим знаменателем.
Характеристический полином --- это \textbf{числитель} суммы.

Для типовой задачи
\[
W_1(s)=\frac{1}{s^2+a_1 s+a_2},\qquad
W_2(s)=K_1+\frac{K_2}{s}=\frac{K_1 s+K_2}{s}
\]
сразу получается
\[
P(s)=s^3+a_1 s^2+(a_2+K_1)s+K_2.
\]

\begin{itemize}
  \item Не работай с \(W(s)\) ``на глаз'' по схеме. Всегда своди задачу к полиному
        \(P(s)=\alpha_n s^n+\cdots+\alpha_0\).
  \item Михайлов и Харитонов применяются именно к \textbf{замкнутой} системе, то есть к \(P(s)\).
  \item Член \(K_2/s\) поднимает степень: не забудь умножить на знаменатель \(W_2\).
\end{itemize}

%---------------------------------------------------------------------
\section{Классический критерий Михайлова}

Подставь \(s=\jomega\):
\[
P(\jomega)=U(\omega)+jV(\omega).
\]

\subsection{Как получить \texorpdfstring{\(U\) и \(V\)}{U и V}}

Чётные степени \(s^{2k}\) идут в \(U\) со знаком \((-1)^k\);
нечётные \(s^{2k+1}\) идут в \(V\) со знаком \((-1)^k\) и множителем \(\omega\).

Для \(P(s)=s^3+\alpha s^2+\beta s+\gamma\):
\[
U(\omega)=\gamma-\alpha\omega^2,\qquad
V(\omega)=\omega(\beta-\omega^2).
\]

Для \(P(s)=s^4+\alpha s^3+\beta s^2+\gamma s+\delta\):
\[
U(\omega)=\delta-\beta\omega^2+\omega^4,\qquad
V(\omega)=\omega(\gamma-\alpha\omega^2).
\]

\subsection{Формулировка критерия}

Система асимптотически устойчива, если при \(\omega:0\to+\infty\) годограф \(P(\jomega)\):
\begin{enumerate}
  \item начинается на \textbf{положительной вещественной оси}: \(U(0)>0\), \(V(0)=0\);
  \item проходит \textbf{последовательно} \(n\) квадрантов против часовой стрелки:
        \(\mathrm{I}\to\mathrm{II}\to\mathrm{III}\to\cdots\);
  \item \textbf{не проходит через начало координат} (иначе есть корень на оси \(j\omega\));
  \item уходит в бесконечность в нужном квадранте.
\end{enumerate}

Эквивалентно: \(\arg P(\jomega)\) монотонно растёт на \(\dfrac{n\pi}{2}\).

\subsection{Типичные отказы годографа}

\noindent
\renewcommand{\arraystretch}{1.25}%
\begin{tabularx}{\textwidth}{>{\raggedright\arraybackslash}X >{\raggedright\arraybackslash}X}
\toprule
\textbf{Что видно на годографе} & \textbf{Вывод} \\
\midrule
Проходит через \(0\) & Граница устойчивости \\
Пропускает квадрант / идёт по часовой & Неустойчивость \\
Стартует не с \(U(0)>0\) & Не все коэффициенты положительны \\
Возвращается в уже пройденный квадрант & Неустойчивость \\
На границе \(\alpha\beta=\gamma\) (для \(n=3\)) годограф идёт в \(0\) & Нет асимптотической устойчивости \\
\bottomrule
\end{tabularx}

\paragraph{Как искать пересечения осей.}
\begin{itemize}
  \item \(U(\omega)=0\), \(V(\omega)>0\) --- переход \(\mathrm{I}\to\mathrm{II}\);
  \item \(V(\omega)=0\), \(U(\omega)<0\) --- переход \(\mathrm{II}\to\mathrm{III}\).
\end{itemize}
Если порядок пересечений ломается --- система неустойчива.

%---------------------------------------------------------------------
\section{Связь с критерием Гурвица}

Для низкой степени Михайлов почти всегда сводится к неравенствам Гурвица.
Их полезно знать наизусть и использовать как быструю проверку.

\paragraph{Степень 2.}
\(s^2+\alpha s+\beta\) устойчив \(\Leftrightarrow\) \(\alpha>0,\ \beta>0\).

\paragraph{Степень 3.}
\(s^3+\alpha s^2+\beta s+\gamma\) устойчив \(\Leftrightarrow\)
\[
\alpha>0,\quad \beta>0,\quad \gamma>0,\quad \alpha\beta>\gamma.
\]
На годографе условие \(\alpha\beta>\gamma\) означает: \(U\) меняет знак \textbf{раньше}, чем \(V\),
то есть переход \(\mathrm{I}\to\mathrm{II}\) происходит при ещё положительном \(V\).

\paragraph{Степень 4.}
\(s^4+\alpha s^3+\beta s^2+\gamma s+\delta\) устойчив \(\Leftrightarrow\)
\[
\alpha,\beta,\gamma,\delta>0,\qquad
\alpha\beta>\gamma,\qquad
\alpha\beta\gamma-\alpha^2\delta-\gamma^2>0.
\]

\paragraph{Правило.}
Для \(n\le 3\) сначала выпиши неравенства, потом подтверди их картинкой Михайлова.
Для отчёта нужна именно картина годографа.
Положительность всех коэффициентов --- только необходимое условие; для \(n\ge 3\) этого мало.

%---------------------------------------------------------------------
\section{Робастная постановка: интервальный полином}

Параметры \(a_i,K_i\) заданы интервалами. Тогда коэффициенты \(P(s)\) тоже лежат в интервалах:
\[
P(s)=\sum_{k=0}^{n}\alpha_k s^k,\qquad
\alpha_k\in[\underline{\alpha}_k,\overline{\alpha}_k].
\]

\subsection{Как пересчитывать интервалы коэффициентов}

\begin{itemize}
  \item Если коэффициент равен параметру, интервал тот же:
        \(\alpha=a_1\in[1,2]\).
  \item Если сумма независимых слагаемых:
        \(\min(a_2+K_1)=\min a_2+\min K_1\),
        \(\max(a_2+K_1)=\max a_2+\max K_1\).
  \item Если произведение положительных:
        \(\min(ab)=\min a\cdot\min b\),
        \(\max(ab)=\max a\cdot\max b\).
  \item Если параметр фиксирован (например \(K_2\in[3,3]\)), коэффициент --- \textbf{константа}.
\end{itemize}

\subsection{Независимость коэффициентов}

Если разные коэффициенты зависят от \textbf{разных} параметров, они независимы,
и можно брать все углы прямоугольника коэффициентов.

Если один параметр входит в два коэффициента, не все углы достижимы.
Тогда теорема Харитонова может быть консервативна: неустойчивый харитоновский полином
может не соответствовать ни одному реальному набору параметров.
В этом случае проверяй вершины и рёбра по \textbf{исходным} параметрам.

%---------------------------------------------------------------------
\section{Теорема Харитонова}

Интервальный полином робастно устойчив тогда и только тогда, когда устойчивы
\textbf{четыре} полинома Харитонова:
\begin{align*}
K_1(s)&=\underline{\alpha}_0+\underline{\alpha}_1 s+\overline{\alpha}_2 s^2+\overline{\alpha}_3 s^3+\underline{\alpha}_4 s^4+\cdots,\\
K_2(s)&=\underline{\alpha}_0+\overline{\alpha}_1 s+\overline{\alpha}_2 s^2+\underline{\alpha}_3 s^3+\underline{\alpha}_4 s^4+\cdots,\\
K_3(s)&=\overline{\alpha}_0+\underline{\alpha}_1 s+\underline{\alpha}_2 s^2+\overline{\alpha}_3 s^3+\overline{\alpha}_4 s^4+\cdots,\\
K_4(s)&=\overline{\alpha}_0+\overline{\alpha}_1 s+\underline{\alpha}_2 s^2+\underline{\alpha}_3 s^3+\overline{\alpha}_4 s^4+\cdots.
\end{align*}

\paragraph{Мнемоника.}
Коэффициенты чередуются парами ``мин--мин / макс--макс'', со сдвигом на один индекс.

Каждый \(K_i\) проверяй критерием Михайлова (или Гурвица).

\subsection{Как применять}

\begin{enumerate}
  \item Выпиши четыре полинома.
  \item Если \textbf{хотя бы один} неустойчив --- семейство \textbf{не} робастно устойчиво.
        Контрпример готов.
  \item Если все четыре устойчивы \textbf{и} коэффициенты независимы ---
        семейство робастно устойчиво.
  \item Если коэффициенты зависимы --- Харитонов достаточен, но не всегда необходим.
\end{enumerate}

Для \(n=3\) часто хватает одного неравенства \(\alpha\beta>\gamma\) на самом ``плохом'' углу:
сравни \(\min(\alpha\beta)=\underline\alpha\cdot\underline\beta\) с \(\overline\gamma\).

%---------------------------------------------------------------------
\section{Михайлов для семейства полиномов}

Годографы всех полиномов семейства заполняют ``трубку''.
Робастная устойчивость есть, если:
\begin{itemize}
  \item ни один годограф не проходит через \(0\);
  \item все годографы идут по тем же \(n\) квадрантам, что и номинальный устойчивый.
\end{itemize}

На практике проверяют граничные (экстремальные) полиномы: углы интервального прямоугольника.
Для независимых коэффициентов это как раз харитоновские полиномы.

%---------------------------------------------------------------------
\section{Алгоритм решения задачи с доски}

\begin{enumerate}
  \item По схеме выписать \(W_1(s)\) и \(W_2(s)\).
  \item Найти \(P(s)=1+W_1 W_2\), привести к полиному.
  \item Записать интервалы коэффициентов \(\alpha_k\).
  \item Проверить необходимое условие: все \(\alpha_k>0\) на всём семействе.
        Если какой-то коэффициент может быть \(\le 0\) --- сразу неустойчивость.
  \item Составить четыре полинома Харитонова.
  \item Для каждого построить \(U(\omega)\), \(V(\omega)\) и проверить порядок квадрантов
        (или применить Гурвиц).
  \item Если всё устойчиво --- система \textbf{робастно устойчива}.
        Если нет --- указать конкретный набор \(a_i,K_i\) (контрпример)
        и показать, как ломается годограф Михайлова.
  \item По желанию выписать область устойчивости внутри бокса параметров,
        например \(a_1(a_2+K_1)>K_2\).
\end{enumerate}

%---------------------------------------------------------------------
\section{Разобранный пример}

Дано:
\[
W_1(s)=\frac{1}{s^2+a_1 s+a_2},\qquad
W_2(s)=K_1+\frac{K_2}{s},
\]
\[
a_1\in[1,2],\quad
a_2\in[2,3],\quad
K_1\in[0;2{,}5],\quad
K_2=3.
\]

\subsection{Характеристический полином}
\[
P(s)=s^3+a_1 s^2+(a_2+K_1)s+3.
\]
Интервалы коэффициентов:
\[
\alpha=a_1\in[1,2],\qquad
\beta=a_2+K_1\in[2;5{,}5],\qquad
\gamma=3.
\]

\subsection{Условие Михайлова / Гурвица}
\[
a_1>0,\quad a_2+K_1>0,\quad a_1(a_2+K_1)>3.
\]

\subsection{Робастная проверка}
\[
\min a_1(a_2+K_1)=1\cdot 2=2<3.
\]
Семейство \textbf{не} робастно устойчиво.

\paragraph{Контрпример:}
\(a_1=1\), \(a_2=2\), \(K_1=0\), \(K_2=3\):
\[
P(s)=s^3+s^2+2s+3,\qquad
U=3-\omega^2,\qquad
V=\omega(2-\omega^2).
\]
При \(\omega=\sqrt{2}\) имеем \(V=0\), \(U=1>0\): возврат на положительную вещественную ось
и уход в IV квадрант вместо II. По Михайлову система неустойчива.

По Харитонову один из четырёх полиномов --- как раз \(s^3+s^2+2s+3\) --- неустойчив,
поэтому семейство не робастно устойчиво.

\subsection{Область устойчивости внутри бокса}
Система устойчива только при \(a_1(a_2+K_1)>3\).
Например, при \(a_1=2\) --- на всём \([a_2,K_1]\);
при \(a_1=1\) --- лишь при \(a_2+K_1>3\).

%---------------------------------------------------------------------
\section{Частые ловушки}

\begin{itemize}
  \item Забыть умножить на знаменатель \(W_2\) (член \(K_2/s\) даёт лишний множитель \(s\)).
  \item Считать \(K_2\in[3,3]\) интервалом --- это константа.
  \item Верно брать \(\min(a_2+K_1)=\min a_2+\min K_1\), но забыть, что
        \(\max=\max a_2+\max K_1\).
  \item Путать устойчивость разомкнутой и замкнутой систем.
        Михайлов --- про замкнутую, то есть про \(P(s)\).
  \item Считать, что ``все коэффициенты положительны'' означает устойчивость.
        Для \(n\ge 3\) нужно ещё условие \(\alpha\beta>\gamma\).
  \item На границе \(\alpha\beta=\gamma\) годограф проходит через \(0\) ---
        это \textbf{не} асимптотическая устойчивость.
\end{itemize}

\vspace{1.2em}
\noindent\textcolor{ruleblue}{\rule{\textwidth}{0.6pt}}\\[0.4em]
{\small\color{rulegray}
Краткая шпаргалка. Для задачи с доски достаточно: \(P(s)\) \(\to\) интервалы \(\alpha_k\) \(\to\)
четыре полинома Харитонова \(\to\) годограф Михайлова / неравенства Гурвица \(\to\)
вердикт о робастной устойчивости и контрпример при отказе.}

\end{document}
